\section{Experimental Setup}
\label{sec:setup}

\paragraph{Tasks.}
We use 20 short computer tasks of two kinds: 9 coding tasks (for example, fix a bug so the tests
pass, implement a function, rename a function everywhere it is used) and 11 PC tasks on files,
shell and data (for example, sort files into folders by type, count matching lines in a log, sum a
column of a CSV file). Each task comes with starting files written into a fresh folder, and an
automatic checker that looks only at the final state of the folder and returns pass or fail. Where
a task gives test files, the checker also confirms they were not changed. Every checker was tested
without any model: it must fail on the unchanged starting folder and pass after a reference
solution. The tasks are split into
15 train tasks (7 coding, 8 PC), which the harness may learn from, and 5 held-out tasks (2 coding,
3 PC), which are only evaluated and never shown to the refiner. The task set was fixed before any
results were seen. Appendix~\ref{app:tasks} lists all tasks.

\paragraph{Protocol.}
Each condition is run with 3 seeds. For each seed:
\begin{enumerate}
  \item \textbf{Learn.} The agent runs the 15 train tasks once each, in a shuffled order (the same
        three orders for every condition and model). The harness is edited between tasks according
        to the condition.
  \item \textbf{Evaluate.} The final harness is frozen, and the agent runs all 20 tasks once.
\end{enumerate}
The static baseline C0 has no learning phase and runs the 20 tasks once per seed. All reported
success rates are from the evaluation phase, giving 60 task runs (45 train, 15 held-out) per
condition.

\paragraph{Models.}
We use two small worker models, \texttt{qwen3.5:0.8b-mlx} (0.8B parameters) and
\texttt{gemma4:e2b-it-qat}, run locally with Ollama~\cite{ollama}. The separate refiner is
\texttt{gemma4:e4b} for both workers, and the memory database uses \texttt{nomic-embed-text}.
All model calls use temperature 0.7, a fixed context size of 32{,}768 tokens, at most 4{,}096
output tokens per reply, and the run's seed as the sampling seed. Thinking mode is off. We tested all conditions on \texttt{qwen3.5:0.8b-mlx}; on \texttt{gemma4:e2b-it-qat} we ran
C0, C1 and C8.

\paragraph{Hardware and serving.}
All runs use consumer hardware. Every condition of one worker model runs on the same Ollama server
with the same model file (we record the server version and model digest of every run), because the
same seed can sample differently on a different server. For the gemma worker, the refiner and
embedding model run on a second Ollama server, the same one used for the qwen runs. Each model is
unloaded from memory after every task and every refiner call, to keep memory use flat over a run.
A run that stops because of an infrastructure problem (lost connection, machine sleep) is never
counted: it is set aside and the whole condition is rerun from the start.

\paragraph{Metrics.}
\begin{itemize}
  \item \textbf{Success rate} on all 20 tasks, on train tasks and on held-out tasks, with 95\%
        Wilson score intervals~\cite{wilson1927probable}.
  \item \textbf{Paired comparison with C0.} Each (task, seed) pair is matched across conditions.
        We count the pairs where a condition passes and C0 fails (better) and the opposite (worse),
        and test the difference with an exact two-sided McNemar test~\cite{mcnemar1947note}.
  \item \textbf{Forgetting.} The number of train tasks that passed during the learning phase but
        failed at evaluation with the final harness.
  \item \textbf{Cost.} Mean wall-clock seconds and prompt tokens per evaluation task, and minutes
        of learning (including refiner calls) per seed.
\end{itemize}
